\documentclass[11pt]{article}%
\usepackage{amsfonts}
\usepackage{amssymb}
\usepackage{graphicx}
\usepackage{setspace}
\usepackage[round]{natbib}
\usepackage{amsmath}%
\usepackage{amsthm}%
\usepackage{palatino}
\usepackage{paralist}
\usepackage{multirow}
\usepackage{booktabs}
\usepackage{dsfont}
\usepackage{indentfirst}
\usepackage{enumerate}
\usepackage{longtable}
\usepackage{caption}
\usepackage{subcaption}
\usepackage{mwe}
\usepackage[hyperindex,breaklinks]{hyperref}
\hypersetup{colorlinks=true,       % false: boxed links; true: colored links
   linkcolor=red,       
    citecolor=blue,        % color of links to bibliography
    filecolor=magenta,      % color of file links
    urlcolor=cyan           % color of external links
}  

\newcommand{\E}{\text{E}}
\newcommand{\I}{\text{I}}
\newcommand{\quantile}{\text{Quantile}}
\newcommand{\defeq}{\mathrel{\mathop:}=}
\newcommand{\norm}[1]{\left\lVert#1\right\rVert}

\usepackage{setspace}
\doublespacing

\bibliographystyle{chicago}

\setcounter{MaxMatrixCols}{10}
\textwidth=6.6in
\textheight=8.9in
\headheight=0.0in
\oddsidemargin=0.0in
\headsep=0.0in
\topmargin=0.0in
\setlength\parindent{24pt}

\setlength{\abovedisplayskip}{5pt}
\setlength{\belowdisplayskip}{5pt}

\title{Bayesian Analysis of Exponential Reciprocal Gamma Models}

\author{ 
}

\date{\today}

\begin{document}
	
\maketitle
%%%%%%%%%%%%%%%%%%%%%
\begin{abstract}
%%%%%%%%%%%%%%%%%%%%%
\noindent Blabla


\bigskip
\noindent {\bf Key Words:} 
\end{abstract}

\newpage
\section{Introduction}
Estimating parameters in a Gamma function \cite{blei2003latent} \cite{escobar1995bayesian}

\cite{bondesson1982simulation} \cite{bondesson1992generalized} \cite{walker2000miscellanea} \cite{ferguson1972representation} \cite{windle2014sampling} \cite{polson2011data} \cite{minka2000estimating}

\cite{roynette2005couples} \cite{hartman1976completely} \cite{polson2013bayesian} 

Generalized Gamma convolution (GGC) can be approximated by a compounded Poisson process. 

\section{Exponential Reciprocal Gamma Distribution}

\begin{equation}
E[e^{-sX}] = \psi(s) = \exp\left\{-as + \int_{0}^\infty (e^{-sy} - 1) N(dy)\right\}
\end{equation}


By Malmsten's formula 
\begin{equation}
\begin{aligned}
\frac{1}{\Gamma(a + s)} &= \exp\left\{\int_{0}^{\infty} (e^{-sx} + \lambda x - 1) \frac{e^{-ax}}{x}dx \right\}\\
&= \int_{0}^{\infty} (1 - e^{-\frac{1}{2} s^2 t}) \frac{1 - e^{-\frac{1}{2} \frac{x^2}{t}}}{\sqrt{2\pi t^3}}dt
\end{aligned}
\end{equation}
Let $\frac{1 - e^{-\frac{1}{2} \frac{x^2}{t}}}{\sqrt{2\pi t^3}}dt = U(dt)$


\begin{equation}
\frac{1 - e^{-\frac{1}{2}x^2/t}}{\sqrt{t}} = \int_0^{\infty} e^{-tz}\frac{2\sin^2(x \sqrt{z / 2})}{\sqrt{\pi z}} dz
\end{equation}
 
\section{Latent Variable Representation of Dirichlet Multinomial Model}
Posterior
\begin{equation}
P(\mathbf{x} \mid \boldsymbol{\alpha}) = \frac{\Gamma(\sum_{k=1}^K \alpha_k)}{\Gamma(N + \sum_{k=1}^K \alpha_k)} \prod_{k=1}^K \frac{\Gamma(N_k + \alpha_k)}{\Gamma(\alpha_k)} e^{-\gamma \sum_{k=1}^K \alpha_k}
\end{equation}



\section{Simulations}

\section{Application}

\section{Discussion}

\newpage
\bibliography{ERgamma} 
\end{document}



